\section{Prime dilation and nonvanishing}\label{sec:dilation}

The previous section makes the forms small after their denominators
are cleared. To obtain a contradiction, we must also ensure that
the eventual rational combination is nonzero infinitely often.
At indices divisible by a large prime, the constant coefficient
and the beta coefficients have different valuations. The seed
constant survives a common scaling while the beta coefficients
vanish modulo the prime. This separation is the principle behind
the auxiliary-prime criterion of Lai and Sprang~\cite[Lemma~2.3]{LS25}.
We obtain it by extending the prime-dilation calculation
of~\cite[Section~6]{Tcount} to the differentiated functions.

Write $\Zp=\{a/b\in\Q:p\nmid b\}$ and let $v_p$ be the
usual prime valuation. Recall $\tau_0=S-m-E$ and
$\tau=\tau_0+k=S-1-E>k$ from \eqref{eq:taus}.

\begin{proposition}\label{prop:dilation}
Fix a positive odd integer $\ell$. For all sufficiently large
primes $p$,
\begin{align}
 p^\tau\rho_\nu(\ell p)&\in p^\nu\Zp,
                                      \label{eq:rho-dilation}\\
 p^\tau q_b(\ell p)&\equiv
     \ch(p)(\ell+1)^{-S}q_b(\ell)\pmod p.
                                      \label{eq:q-dilation}
\end{align}
\end{proposition}

\subsection{The undifferentiated function}

Take $p$ large enough in terms of the fixed parameters and $\ell$;
in particular, it divides none of their denominators and
$p>\ell+1$. Both $(\ell p)!$ and $(\ell p+1)!$ have
$p$-valuation $\ell$. After removing $p^\ell$, their units
reduce to $(-1)^\ell\ell!$, by Wilson's theorem.
Fermat's theorem treats the fixed powers of $4$ and $K$.
Thus $C_{\ell p}$ has valuation $(S-mDL)\ell$, and its unit is
congruent to
\begin{equation}\label{eq:C-unit}
 (-1)^{\ell(S-mDL)}\frac{C_\ell}{(\ell+1)^S}.
\end{equation}

In a numerator block of fraction class $v/d$, a factor of
$P_{\ell p}(pz)$ is selected by $p$ when $dj+v\equiv0\pmod p$.
After division by $p$, its fraction class becomes $v'/d$, where
$pv'\equiv v\pmod d$. This permutes the reduced classes for
each denominator, and hence permutes $\Theta$.
Each block has $L\ell$ selected factors. Together with the
extra factor $p(z-r\ell)$, they give
$p^{DL\ell+1}P_\ell(z)$. The unselected numerator units
contribute $(-1)^{DL\ell}$, since each complete block of nonzero
residues has product $-1$. For $P_{\ell p}^m$, raise both
contributions to the $m$th power.

The selected denominator indices are $h=py$, $0\le y\le\ell$.
They give
\[
 p^{S(\ell+1)}\prod_{y=0}^\ell(z+4y)^S,
\]
and the unselected denominator unit is $(-1)^{S\ell}$.
The signs cancel, while the net valuation is
\[
 (S-mDL)\ell+m(DL\ell+1)-S(\ell+1)=m-S.
\]
Using homogeneity \eqref{eq:homogeneity}, we obtain
\begin{equation}\label{eq:T-dilation}
 p^{\tau_0}T_{\ell p}(pz)
                       \equiv(\ell+1)^{-S}T_\ell(z)\pmod p.
\end{equation}

\subsection{Local coefficient extraction}

To extract coefficients from \eqref{eq:T-dilation}, we need
reduction modulo $p$ to commute with the local series expansion.
At a selected pole $h=py$, put $z=-4y+u$. The factorization above
gives
\begin{equation}\label{eq:local-p-integral}
 u^S p^{\tau_0}T_{\ell p}(p(u-4y))\in\Zp[[u]].
\end{equation}
After the displayed powers are removed, every other denominator
factor has a unit constant term. Its inverse has $p$-integral
series coefficients, and reduction commutes with Taylor extraction.
The coefficient of $u^{S-i}$ in \eqref{eq:local-p-integral} is
$p^{\tau_0-i}a_{i,py}(\ell p)$, so
\begin{equation}\label{eq:a-dilation}
 p^{\tau_0-i}a_{i,py}(\ell p)
                   \equiv(\ell+1)^{-S}a_{i,y}(\ell)\pmod p.
\end{equation}
Now use the integer binomial coefficient in
\eqref{eq:derivative-coefficients}. With $\nu=i+k$ and
$\tau=\tau_0+k$, this gives
\begin{equation}\label{eq:selected-coefficients}
 p^{\tau-\nu}A_{\nu,py}(\ell p)
                   \equiv(\ell+1)^{-S}A_{\nu,y}(\ell)\pmod p.
\end{equation}
Thus the passage to differentiated coefficients is made after
the justified local extraction.

Consider next a pole $h=x+py$ with $1\le x<p$, and substitute
$t=pz-4x$. The extra numerator factor is a unit, the numerator
blocks have $DL\ell$ selected factors, and the denominator has
$\ell$ selected factors, each of order $S$. Including $C_{\ell p}$,
the net valuation is zero. Local expansion of $R_{\ell p}$ at
$-4h$ shows that its partial-fraction coefficient of order $i$
has valuation at least $i$.
Multiplication by $W_{\ell p}$ preserves this inequality: its
$j$th Taylor coefficient is an integer and multiplies a
coefficient of $R_{\ell p}$ of order $i+j$. Differentiation then gives
\begin{equation}\label{eq:other-coefficients}
 v_p\bigl(A_{\nu,h}(\ell p)\bigr)\ge\nu-k
                                  \qquad(p\nmid h).
\end{equation}
This weaker bound suffices because $\tau>k$. Combining
\eqref{eq:selected-coefficients} and \eqref{eq:other-coefficients}
over all poles proves \eqref{eq:rho-dilation}.

\subsection{The surviving rational constant}

Use \eqref{eq:q}, taking $p>4B\ell$ as well. Then every
$j\le4b\ell p$ is less than $p^2$, so $v_p(j)\le1$.
At a nonselected pole, multiplication by $p^\tau$ leaves
valuation at least $\tau-k>0$, even if $p\mid j$.
At a selected pole, a term with $p\nmid j$ has valuation
at least $\nu>0$, by \eqref{eq:selected-coefficients}.

Only the terms with $h=py$ and $j=pj'$ survive. Their upper
limit becomes $j'\le4by$, and
$\ch(pj')=\ch(p)\ch(j')$. Equation~\eqref{eq:selected-coefficients}
therefore turns the surviving sum into
$\ch(p)(\ell+1)^{-S}$ times the constant formula at index
$\ell$. This is \eqref{eq:q-dilation}, completing the proof
of Proposition~\ref{prop:dilation}.
